\section{Window experiment planning}
\label{app:planning}

The full-access window (expected 2026-10-29 17:00 to 10-30 17:00) is planned so that the most important
results come first and the plan degrades gracefully if the window is cut short. The detailed checklist,
with the run-by-run queue and owners, is in \texttt{scripts/analysis/EXPERIMENTS.md}.

\paragraph{Prioritization.} The campaign started as a validation of two new milabench workloads (Qwen3.8-27B
on TorchTitan and GLM-5 on Megatron) at a scale they had never run at. Full access to Rorqual turned it into
an opportunity to also use these workloads as probes of the questions that matter to the cluster's
administrators. The ranking keeps both purposes in view, with four criteria: whether a result shows that
each workload runs and measures correctly at full scale (the original purpose), how directly it answers the
cluster questions (does Rorqual's network limit large-scale training, and what follows for allocation and
for the exascale evaluation), whether it informs a decision, and whether only the window can produce its
data. The first criterion does not change the order: the tier-1 runs already include each workload on every
healthy node. Two rules refine the ranking:
\begin{itemize}[leftmargin=*, itemsep=1pt]
  \item \textbf{Tier 0 before anything else}: the health check of every node is a prerequisite, not a
        result. It sets the healthy set, the exclusions and the node lists of every later run, and it is
        the first snapshot of the state of the cluster.
  \item \textbf{GLM-5 first when near-equal}: GLM-5 runs are more complex and more likely to fail; running
        them first leaves more time to fix them before the window ends.
\end{itemize}

\begin{table}[h]
  \caption{Figure priority. Tier 1 alone answers the main questions; tier 2 explains them; tier 3 is
    context, or already covered by preliminary data.}
  \label{tab:priority}
  \centering\small
  \begin{tabular}{cll}
    \toprule
    Tier & Figure & Role \\
    \midrule
    0 & Window health check (data) & sanity check, state of the cluster, node lists \\
    \midrule
    1 & \Cref{fig:expert-placement} GLM-5 expert placement & using every node vs.\ the largest aligned subset \\
    1 & \Cref{fig:scaling} Weak scaling & throughput per GPU up to the full machine \\
    1 & \Cref{fig:placement-at-scale} At scale: balance or pack & allocation of large jobs \\
    1 & \Cref{fig:placement-scaling} More nodes vs.\ crossing switches & where the scaling loss comes from \\
    \midrule
    2 & \Cref{fig:comm} Bandwidth by scope & physical cause (one HCA per node, spine) \\
    2 & \Cref{fig:health-window} Window health check & validity of the rest \\
    2 & \Cref{fig:moe} GLM-5 scaling, EP size, FP8 & MoE at scale \\
    2 & \Cref{fig:expert-groups} All-to-all per EP group & mechanism behind \cref{fig:expert-placement} \\
    2 & \Cref{fig:decomposition} Step decomposition & communication-bound step \\
    \midrule
    3 & \Cref{fig:batch} Per-GPU batch & answered by preliminary data \\
    3 & \Cref{fig:startup} Startup & fed by every run \\
    3 & \Cref{fig:stragglers} Stragglers & node telemetry from regular runs \\
    3 & \Cref{fig:stability} Stability & conditional on window length \\
    3 & \Cref{fig:health-gpus} Per-GPU compute by cluster & premise shown by preliminary data \\
    3 & \Cref{fig:health-tt} Preliminary health checks & background \\
    \bottomrule
  \end{tabular}
\end{table}

\paragraph{Roles.} Three people run the window: one for health checks and monitoring (runs the health
checks, watches GPU telemetry and logs during every run, regenerates the quick figures after each block and
compares them with expectations, flags anything suspicious), one for the Qwen workload and one for the GLM-5
workload.

\paragraph{Rules.}
\begin{itemize}[leftmargin=*, itemsep=1pt]
  \item \textbf{Small blocks, alternating workloads}: each block is one comparison (10--36 minutes) and
        launches from a prepared command; GLM-5 and Qwen blocks alternate, so one team prepares while the
        other runs.
  \item \textbf{Skip on failure}: a failed or suspicious block is logged and the next block of the
        \emph{other} workload starts; the failed block returns to the queue once fixed.
  \item \textbf{First pass, then repeats}: pass 1 runs every tier-0 and tier-1 comparison once, so every key
        figure has data early; pass 2 repeats it; then the remaining points and the optional runs.
  \item \textbf{Parallel runs}: 1-node runs (smoke tests, re-checks) use idle nodes at any time; a run that
        stays on one switch runs next to any run that does not use that switch; runs that cross the spine run
        one at a time. The assumption that a switch does not slow traffic between its own nodes is checked in
        the window (the same one-switch run with and without a concurrent run elsewhere).
  \item \textbf{Exclusions}: a node flagged by the health check is excluded and the node lists are
        recomputed; a comparison interrupted by an exclusion is rerun on the new node set.
\end{itemize}

\begin{table}[h]
  \caption{Window phases (Qwen runs measured at $\sim$12 minutes; GLM-5 runs assumed at $\sim$15 minutes until
    measured), with parallel runs. About 11.5 hours of execution, leaving $\sim$12.5 hours of the window for
    fixes and failures.}
  \label{tab:phases}
  \centering\small
  \begin{tabular}{lp{0.72\linewidth}r}
    \toprule
    Phase & Content & Time \\
    \midrule
    Pass 1 & health check; GLM-5 aligned, crossed, mixed, every node; Qwen every node, H27/P27, S3/S12/O12, H12,
             leaf-to-leaf controls; whole-machine health checks & 2.5 h \\
    Pass 2 & second repetition of pass 1 (one-switch check run alone); health re-check & 2.4 h \\
    P2 & GLM-5 FP8; Qwen H15/P15, H48/P48, 6- and 9-node placements, O16, O3; their repeats & 4.5 h \\
    P3 & GLM-5 EP16; bridge to past data with gradient accumulation; 1-hour stability run; final health check & 2.1 h \\
    \bottomrule
  \end{tabular}
\end{table}

\paragraph{Coverage.} \Cref{tab:coverage} states, for each figure, what the report would contain if the
window stopped after each phase.

\begin{table}[h]
  \caption{Report coverage by phase (cumulative). \emph{none}: no window data; \emph{usable}: supports a
    statement, with the limit given; \emph{complete}: as designed.}
  \label{tab:coverage}
  \centering\footnotesize
  \setlength{\tabcolsep}{3pt}
  \begin{tabular}{c>{\raggedright}p{3.1cm}>{\raggedright}p{2.35cm}>{\raggedright}p{2.35cm}>{\raggedright}p{2.2cm}>{\raggedright\arraybackslash}p{2.0cm}}
    \toprule
    Tier & Figure & After pass 1 & After pass 2 & After P2 & After P3 \\
    \midrule
    0 & Health check (data) & complete & complete & complete & complete \\
    1 & \ref{fig:expert-placement} Expert placement & usable: 1 run each, no significance & \textbf{complete} & complete & complete \\
    1 & \ref{fig:scaling} Weak scaling & usable: 12--264 GPUs, 5 scales & usable, repeats & \textbf{complete} & complete \\
    1 & \ref{fig:placement-at-scale} Balance or pack & usable: 27 nodes, 1 run & usable: with run-to-run spread & \textbf{complete}: 15, 27, 48 & complete \\
    1 & \ref{fig:placement-scaling} Nodes vs.\ crossing & usable: 3 and 12 nodes, controls & usable, repeats & \textbf{complete}: 6, 9, 16 & complete \\
    2 & \ref{fig:comm} Bandwidth by scope & \textbf{complete} & complete & complete & complete \\
    2 & \ref{fig:health-window} Window health check & \textbf{complete} & complete & complete & complete \\
    2 & \ref{fig:moe} GLM-5 scaling & usable: EP8 BF16 & usable, repeats & usable: + FP8 & \textbf{complete}: + EP16 \\
    2 & \ref{fig:expert-groups} All-to-all per group & complete if timers logged & complete & complete & complete \\
    2 & \ref{fig:decomposition} Step decomposition & usable (estimated) & usable & usable & \textbf{complete} if measured \\
    3 & \ref{fig:batch} Per-GPU batch & usable (preliminary) & same & same & \textbf{complete} \\
    3 & \ref{fig:startup} Startup & usable & usable & \textbf{complete} & complete \\
    3 & \ref{fig:stragglers} Stragglers & usable (5-min windows) & usable & usable & \textbf{complete} \\
    3 & \ref{fig:stability} Stability & none & none & none & \textbf{complete} \\
    3 & \ref{fig:health-gpus} Per-GPU compute & complete if tool reports it & same & same & same \\
    3 & \ref{fig:health-tt} Preliminary health checks & complete & complete & complete & complete \\
    \bottomrule
  \end{tabular}
\end{table}

\paragraph{What is lost if the window stops early.} Before pass 2 ends, the placement results
(\cref{fig:expert-placement,fig:placement-at-scale}) have no significance test: differences of 0.2--1\,\% cannot
be told from run-to-run noise with one run each. Before P2 ends, the tier-1 curves stay sparse and GLM-5 has
no FP8 result. Before P3, there is no stability figure, the step decomposition stays estimated, GLM-5 has no
EP-size comparison at full scale, and the batch figure has no window point. Independently of the window,
\cref{fig:expert-groups} depends on a tool change made before it (per-group all-to-all timers in Megatron).
